% Capítulo 1, Introdução

% \section{Contextualização}
\label{sec-contextualizacao}

Aplicações que servem requisições autenticadas mantêm, para cada identidade, um conjunto de estados de curta duração. \revini{INT-01}Esses estados incluem sessões abertas, contadores de tentativas falhas, metadados de renovação e marcas de bloqueio, e existem para que cada nova requisição possa ser validada como continuação legítima de uma autenticação anterior \citep{nist_sp_800_63b4_2025, owasp_session_mgmt_cheatsheet}.\revfim{INT-01} Em arquiteturas com várias instâncias atrás de um balanceador de carga, esse estado não pode ficar restrito ao processo que o criou, sob pena de afinidade rígida do cliente ao processo ou de perda de sessão em qualquer reinício. A resposta corrente é externalizar o estado para um armazenamento compartilhado com acesso de baixa latência. Os \emph{In-Memory Data Grids} (IMDG) ocupam esse papel: oferecem armazenamento chave-valor em memória, particionamento por chave, replicação configurável e mecanismos de realocação quando a composição do \emph{cluster} muda \citep{infinispan_intro, infinispan_docs_stable}.

A replicação dentro de um IMDG admite dois desenhos. \revini{INT-02}Na replicação síncrona, a confirmação de uma escrita depende da confirmação remota das demais cópias, o que atrela cada operação ao \emph{Round-Trip Time} (RTT) e à variabilidade da rede.\revfim{INT-02} Na replicação assíncrona, ou otimista, a escrita é confirmada localmente e propagada em segundo plano, com ganho de latência percebida e o custo de admitir réplicas que observam o mesmo objeto em estados diferentes. Trabalhos seminais demonstram que essa segunda família é uma escolha consciente de armazenamentos chave-valor de produção, com mecanismos como \emph{vector clocks}, \emph{sloppy quorum} e \emph{hinted handoff} para tolerar falhas sem coordenação síncrona \citep{decandia_dynamo_2007}. Para tratar a reconciliação entre réplicas divergentes, uma alternativa consolidada são os tipos de dados replicados livres de conflito, que garantem convergência determinística do estado quando as operações satisfazem propriedades algébricas adequadas \citep{shapiro_crdt_2011}.

Além do funcionamento sob condições normais, este trabalho considera explicitamente três classes de evento que rompem a estabilidade do \emph{cluster}. A primeira é o \emph{failover}: um nó deixa de responder e o atendimento passa a uma réplica, possivelmente com dados defasados. A segunda é a partição de rede, em que sub-conjuntos de nós deixam de trocar mensagens por um período. A terceira reúne o rebalanceamento e o \emph{state transfer}, eventos planejados ou reativos que alteram o mapeamento entre chaves e nós e podem coexistir com falhas \citep{infinispan_docs_stable}. Cada uma dessas três classes pode afetar tanto propriedades funcionais, relativas ao estado final do sistema, quanto propriedades temporais, relativas à ordem dos eventos observados pelos clientes.  Quando o trabalho exige replicação entre localidades distintas, são considerados os mecanismos documentados para replicação \emph{cross-site} \citep{infinispan_xsite}. 

O problema também tem leitura no eixo da computação sustentável. \revini{INT-18}Consumo de energia e resiliência estão entre os desafios centrais da computação de alto desempenho \citep{navaux_hpc_challenges_2023}, e falhas cobram recursos em reconstruções de estado, repetição de operações e reprocessamento.\revfim{INT-18} O consumo energético desses ambientes está diretamente atrelado à execução de tarefas computacionais. A literatura analisa estratégias de economia de energia em sistemas tolerantes a falhas \citep{ft_hpc_energy_2024}.

Este trabalho aborda essa preocupação por dois caminhos. O primeiro é a redução de retrabalho, ou seja, ao reduzir violações de invariantes e estabilizar o comportamento sob falhas, o sistema evita reexecuções e reprocessamento, e com eles o gasto computacional associado. O segundo é a camada de armazenamento, o estado de sessão é mantido em memória em um \emph{data grid}, e \revini{INT-19}decisões de \emph{cache} têm impacto de carbono mensurável\revfim{INT-19} e podem ser tratadas de forma consciente de carbono \citep{tian_greencache_2026}. O controlador heurístico previsto, que ajusta políticas e parâmetros de configuração em tempo de execução, se alinha a esforços recentes que investigam o autoajuste de parâmetros sob a perspectiva da computação sustentável, com o objetivo de reduzir o custo computacional do processamento \citep{lemos_eradrs_2025}. O objetivo de reduzir retrabalho e estabilizar a cauda de latência é coerente com essa preocupação\revini{INT-27}.\revfim{INT-27}

\section{Problema de pesquisa}
\label{sec-problema}

O problema é formulado em termos de invariantes de domínio do estado de sessão. Um invariante é uma propriedade que deve permanecer verdadeira em todas as execuções permitidas pelo modelo, apesar de concorrência e falhas. Três exemplos guiam o trabalho. A irreversibilidade da invalidação afirma que, depois de um \emph{logout} confirmado, a sessão correspondente não deve retornar ao estado ``válida''. A monotonicidade de contadores afirma que tentativas falhas registradas não devem diminuir como resultado de reconciliações entre réplicas. A consistência de bloqueio afirma que uma identidade marcada como bloqueada não deve ter sessões validadas como ativas enquanto o bloqueio durar. \revini{INT-06}As três propriedades têm respaldo normativo, com tratamento explícito em \citet{nist_sp_800_63b4_2025} e em \citet{owasp_session_mgmt_cheatsheet}, e podem ser violadas quando replicação assíncrona e \emph{failover} interagem.\revfim{INT-06} \revini{INT-07}As três propriedades derivam de exigências de \citet{nist_sp_800_63b4_2025}, como a invalidação dos segredos de sessão no \emph{logout} e a limitação de tentativas falhas, e de \citet{owasp_session_mgmt_cheatsheet}, que exige a invalidação da sessão no servidor no \emph{logout}. A forma replicada dessas propriedades é derivação deste trabalho.\revfim{INT-07}

Diante disso, este trabalho aborda como questão central de pesquisa especificar e verificar formalmente invariantes do estado de sessão em \emph{data grids} sob replicação assíncrona e falhas e, a partir dos contraexemplos produzidos pela verificação e das métricas observadas em experimentos controlados, definir um controlador heurístico que reduza a taxa de violações e estabilize a cauda da distribuição de latência, representada pelos percentis p95 e p99, sem aumentar indevidamente o custo operacional.

A abordagem proposta combina quatro elementos que se encadeiam. O trabalho começa pela especificação formal do núcleo do protocolo de estado de sessão em um sistema de transições com canal assíncrono. Sobre essa especificação, verifica os invariantes declarados por \emph{model checking}, técnica adequada para entrelaçamentos de eventos cuja reprodução por testes é pouco confiável \revini{INT-21}\citep{newcombe_aws_formal_methods_2015}\revfim{INT-21}. A partir dos contraexemplos obtidos, desenha e implementa um mecanismo de decisão heurístico que observa métricas do sistema e escolhe políticas de mitigação.  Por fim, valida o conjunto em um \emph{cluster} Infinispan sob cenários de falha reproduzíveis, com instrumentação que mede a latência, em particular os percentis p95 e p99 \citep{dean_barroso_tail_2013, google_sre_slo}, e que detecta violações de invariantes por verificações no nível da aplicação.

\section{Justificativa}
\label{sec-justificativa}

Externalizar o estado de sessão para um \emph{data grid} replicado faz com que a interação entre concorrência, replicação e falhas se torne um fator determinante da correção observada pelo usuário. \citet{decandia_dynamo_2007} consolidaram, em produção, o uso de replicação assíncrona em armazenamentos chave-valor de larga escala. \revini{INT-22}No Dynamo, uma escrita pode ser confirmada ao cliente antes de chegar a todas as réplicas, bastando a resposta de parte delas, e essa decisão tem como custo réplicas que observam estados defasados e exigem reconciliação posterior.\revfim{INT-22} Quando o dado replicado é o estado de sessão, essa defasagem pode reverter decisões de segurança, como uma invalidação. \citet{shapiro_crdt_2011} mostram que a convergência sem coordenação é possível quando as operações satisfazem propriedades algébricas como comutatividade ou formação de um semirreticulado, o que torna pertinente investigar quais invariantes do estado de sessão admitem garantias análogas e quais exigem mecanismos adicionais.

A pertinência do estudo das falhas que cercam esses invariantes é sustentada por resultados teóricos clássicos e por evidência empírica de produção. No plano teórico, \citet{fischer_lynch_paterson_flp_1985} demonstram que, em um sistema assíncrono no qual processos podem parar silenciosamente, nenhum protocolo determinístico de consenso é totalmente correto, e a prova constrói uma execução que evita a decisão. Na mesma linha, \citet{gilbert_lynch_cap_2002} formalizam o teorema CAP e mostram que, em presença de partição de rede, um sistema distribuído não pode oferecer simultaneamente consistência forte e disponibilidade. \citet{chandra_toueg_failure_detectors_1996} contornam parcialmente a barreira ao introduzir detectores de falha caracterizados por duas propriedades, completude e acurácia, com a classe $\diamond W$ identificada como suficiente para resolver consenso. No plano empírico, \citet{bailis_kingsbury_network_2014} documentam partições, atrasos e perdas de mensagens como eventos recorrentes em infraestruturas de produção, o que reforça que esses cenários devem ser parte do experimento e não exceção dele.

A opção por controlar percentis altos de latência, em vez da média, decorre do impacto desproporcional da cauda da distribuição em serviços compostos por muitos componentes. \revini{INT-09}Um exemplo do livro de engenharia de confiabilidade do Google mostra o problema: a latência típica é da ordem de 50 milissegundos, mas 5\% das requisições são vinte vezes mais lentas, e uma média não acusa a variação na cauda \citep{google_sre_slo}.\revfim{INT-09} \revini{INT-23}\citet{dean_barroso_tail_2013} mostram que o efeito cresce com o número de componentes: se cada servidor responde tipicamente em 10 milissegundos, mas tem p99 de 1 segundo, uma requisição que espera 100 desses servidores passa de 1 segundo em 63\% dos casos.\revfim{INT-23} A formulação dessas metas como objetivos de nível de serviço, com indicadores correspondentes e \emph{error budgets} associados, é prática consolidada de engenharia de confiabilidade \citep{google_sre_slo}. O uso de \emph{model checking}, por sua vez, justifica-se pela capacidade de explorar sistematicamente entrelaçamentos de eventos e produzir contraexemplos concretos para violações, classe de defeitos difícil de reproduzir apenas por testes \revini{INT-24}\citep{newcombe_aws_formal_methods_2015}\revfim{INT-24}. Há registros de aplicação industrial dessa técnica a protocolos de replicação em sistemas distribuídos de larga escala, com defeitos sutis identificados antes do código entrar em produção \revini{INT-25}\citep{newcombe_aws_tla_2014, newcombe_aws_formal_methods_2015}\revfim{INT-25}.

O domínio de aplicação reforça a importância dos invariantes escolhidos. A gerência de sessão em arquiteturas distribuídas em nuvem\revini{INT-26}\revfim{INT-26} é cercada de orientações normativas que descrevem propriedades operacionais do ciclo de vida da sessão, com tratamento detalhado em \revini{INT-10}\citet{nist_sp_800_63b4_2025}\revfim{INT-10} e em \citet{owasp_session_mgmt_cheatsheet}. As três propriedades alvo deste trabalho, irreversibilidade de \emph{logout}, monotonicidade de contadores e consistência de bloqueio, têm respaldo direto nessas fontes normativas\revini{INT-11}, e a formulação delas para o estado replicado é derivação do trabalho\revfim{INT-11}. A combinação proposta, em um mesmo trabalho, de especificação formal, verificação por \emph{model checking}, implementação experimental em um \emph{data grid} de uso corrente e um controlador heurístico orientado a invariantes e a SLO constitui a contribuição pretendida. O posicionamento dessa combinação em relação aos trabalhos correlatos aparece no Capítulo~\ref{cap-metodologia}, a partir do mapeamento sistemático ali descrito.

\section{Objetivos}
\label{sec-objetivos}

\subsection*{Objetivo geral}
\label{sec-objetivo-geral}

Especificar e verificar invariantes do estado de sessão sob replicação assíncrona e \emph{failover} em \emph{data grids}, e avaliar um mecanismo de decisão heurístico que modifique políticas ou ações operacionais para reduzir a ocorrência de violações desses invariantes e controlar métricas de latência sob falhas.

\subsection*{Objetivos específicos}
\label{sec-objetivos-especificos}

Constituem objetivos específicos deste trabalho:

\begin{enumerate}[label=(\arabic*)]
    \item definir um modelo formal do estado de sessão e das operações relevantes (\emph{login}, validação, renovação, invalidação e atualização de contadores e marcadores), incluindo atrasos e reordenação de propagação de atualizações e eventos de falha ou recuperação;
    \item declarar invariantes de segurança, correção e propriedades de convergência, e verificar essas propriedades via \emph{model checking}, registrando contraexemplos quando existirem;
    \item estabelecer uma arquitetura experimental com Infinispan em \emph{cluster}, com suporte a alterações de topologia e cenários de replicação, instrumentando um programa de carga que gera operações concorrentes e detecta violações de invariantes;
    \item implementar um controlador heurístico que, com base em taxas e parâmetros observáveis, selecione entre políticas e aplique ações de mitigação;
    \item comparar experimentalmente o modelo estático com o modelo acrescido do controlador em cenários com injeção de falhas, quantificando taxa de violações, tempo de recuperação e convergência, e estatísticas de latência.
\end{enumerate}



\section{Organização do texto}
\label{sec-organizacao}

O restante deste trabalho está organizado em quatro capítulos. O Capítulo~\ref{cap-fundamentacao} apresenta a fundamentação teórica. O Capítulo~\ref{cap-metodologia} detalha a metodologia, organizada em duas linhas de trabalho, sintetiza os trabalhos correlatos derivados do mapeamento sistemático e apresenta o cronograma das duas etapas do TCC. O Capítulo~\ref{cap-desenvolvimento} reúne o desenvolvimento e os resultados parciais correspondentes ao Trabalho de Conclusão de Curso~I, em particular a arquitetura experimental documentada e a especificação formal inicial. O Capítulo~\ref{cap-conclusao} sintetiza as considerações finais parciais, registra as limitações conhecidas e o planejamento para o Trabalho de Conclusão de Curso~II, cujo plano expandido aparece no Apêndice~\ref{ap-plano-tcc2}. 
